\begin{figure}[H]
    \centering
    \begin{tikzpicture}[x=1.0cm, y=1.0cm]

        % (a) sorted block fits the tile: layers fuse into one pass
        \node[panel] at (-1.05,1.30) {(α) $k' \le$ πλακίδιο};

        \node[blkdrm, minimum width=7.20cm, minimum height=0.42cm] (g1) at (3.60,0.72)
            {\tiny καθολική μνήμη};

        \node[blkcac, minimum width=1.62cm, minimum height=1.05cm] (t1a) at (0.90,-0.55) {};
        \node[blkcac, minimum width=1.62cm, minimum height=1.05cm] (t1b) at (2.70,-0.55) {};
        \node[blkcac, minimum width=1.62cm, minimum height=1.05cm] (t1c) at (4.50,-0.55) {};
        \node[blkcac, minimum width=1.62cm, minimum height=1.05cm] (t1d) at (6.30,-0.55) {};

        \foreach \n in {t1a,t1b,t1c,t1d} {
            \node[notemute, font=\tiny] at (\n) {πλακίδιο};
        }

        \draw[flow] (0.90,0.50) -- (0.90,-0.01);
        \draw[flow] (2.70,0.50) -- (2.70,-0.01);
        \draw[flow] (4.50,0.50) -- (4.50,-0.01);
        \draw[flow] (6.30,0.50) -- (6.30,-0.01);

        \node[note, font=\scriptsize, anchor=west] at (7.55,-0.55)
            {ένα πέρασμα:\\ όλα τα επίπεδα\\ συγχωνεύονται};

        \draw[bracemirror] (0.10,-1.30) -- (1.70,-1.30);
        \node[notemute, font=\tiny, anchor=north] at (0.90,-1.52) {$k'$};

        \begin{scope}[yshift=-3.35cm]
            % (b) block exceeds the tile: distant layers spill to global memory
            \node[panel] at (-1.05,1.30) {(β) $k' >$ πλακίδιο};

            \node[blkdrm, minimum width=7.20cm, minimum height=0.42cm] at (3.60,0.72)
                {\tiny καθολική μνήμη};

            \node[blkcac, minimum width=1.62cm, minimum height=1.05cm] at (0.90,-0.55) {};
            \node[blkcac, minimum width=1.62cm, minimum height=1.05cm] at (2.70,-0.55) {};
            \node[blkcac, minimum width=1.62cm, minimum height=1.05cm] at (4.50,-0.55) {};
            \node[blkcac, minimum width=1.62cm, minimum height=1.05cm] at (6.30,-0.55) {};
            \foreach \x in {0.90,2.70,4.50,6.30} {
                \node[notemute, font=\tiny] at (\x,-0.55) {πλακίδιο};
            }

            \foreach \x in {0.90,2.70,4.50,6.30} {
                \draw[flow] (\x-0.18,0.50) -- (\x-0.18,-0.01);
                \draw[flow] (\x+0.18,-0.01) -- (\x+0.18,0.50);
            }

            \node[note, font=\scriptsize, anchor=west] at (7.55,-0.55)
                {πολλά περάσματα:\\ τα επίπεδα με μεγάλη\\ απόσταση δεν χωρούν};

            \draw[bracemirror] (0.10,-1.30) -- (7.10,-1.30);
            \node[notemute, font=\tiny, anchor=north] at (3.60,-1.52) {$k'$};
        \end{scope}

        \node[note, anchor=north, align=center] at (3.60,-6.05)
            {Η μετρούμενη κίνηση περνά από $7n$ σε $25n$ έως $56n$ όταν
             το $k'$ ξεπεράσει το πλακίδιο.};

    \end{tikzpicture}
    \caption{Η σχέση του ταξινομημένου μπλοκ με το πλακίδιο της κοινής μνήμης. Ένα
    επίπεδο μπορεί να εκτελεστεί χωρίς επιστροφή στην καθολική μνήμη μόνο αν η
    απόσταση των στοιχείων που συγκρίνει χωρά στο πλακίδιο. Όσο τα ταξινομημένα
    μπλοκ μεγέθους $k'$ χωρούν σε αυτό (α), σχεδόν όλα τα επίπεδα συγχωνεύονται σε
    ένα πέρασμα και η κίνηση δεδομένων ταυτίζεται με εκείνη του επεξεργαστή. Όταν
    το $k'$ ξεπεράσει το μέγεθος του πλακιδίου (β), τα επίπεδα με μεγάλη απόσταση
    εκτελούνται με ξεχωριστά περάσματα πάνω από την καθολική μνήμη και η κίνηση
    πολλαπλασιάζεται. Για ακεραίους των 32 \tl{bit} το πλακίδιο είναι 8192
    στοιχεία, οπότε η μετάβαση συμβαίνει μεταξύ των δύο μεγαλύτερων τιμών του
    \tl{k} που εξετάστηκαν.}
    \label{fig:gpu-tiling}
\end{figure}
